% !TEX TS-program = XeLaTeX
% !TEX encoding = UTF-8 Unicode
% Copyright (c) 2014 by Daniel R. Driver. All rights reserved.

\documentclass[titlepage]{article}
% clone the necessary input files: git@github.com:danieldriver/syllabus.git
\newcommand\policy{../policy}
\newcommand\incl{../includes}
\ProvidesFile{variables.tex}[2022/09/11 v2.3 -- Syllabus variables]

\usepackage{xspace} % make manual spaces (like \mycmd\ ) unnecessary
\usepackage{xifthen} % provides \isempty test

% variables for internal use
\newcommand\prof{}
\newcommand\pdegree{}
\newcommand\pphone{}
\newcommand\pemail{}
\newcommand\poffice{}
\newcommand\phours{}
%
\newcommand\ccode{}
\newcommand\ctitle{}
\newcommand\cseries{}
\newcommand\cversion{}
\newcommand\csemester{}
\newcommand\cmeetson{}
\newcommand\cmeetsat{}
\newcommand\cmeetsin{}
\newcommand\cwebsite{}
\newcommand\cdescrip{}
\newcommand\cprereqs{}
\newcommand\edobject{}

% in case of fully online, BTh, or undergraduate courses - https://tex.stackexchange.com/a/5896
\newif\ifonline
\newcommand\Int[2]{\ifonline#1\else#2\fi}
\newif\ifbth
\newcommand\bth[2]{\ifbth#1\else#2\fi}
\newif\ifUG
\newcommand\UG[2]{\ifUG#1\else#2\fi}

% commands for setting variables in the preamble
\newcommand\professor[2][PhD]{
  \renewcommand\pdegree{#1\xspace}
  \renewcommand\prof{#2\xspace}}
\newcommand\phone[1]{
  \renewcommand\pphone{\addfontfeatures{Numbers=Monospaced}#1\xspace}}
\newcommand\email[1]{
  \renewcommand\pemail{\href{mailto:#1}{#1}\xspace}}
\newcommand\officehours[2][AST Library/Teams]{
  \renewcommand\poffice{#1\xspace}
  \renewcommand\phours{\href{https://outlook.office.com/bookwithme/user/b97b147a2e5c43b4a0e200c2cb630b77@astheology.ns.ca?anonymous}{#2}\xspace}%
  }
%
\newcommand\coursecode[2][1.0]{
  \renewcommand\cversion{#1\xspace}
  \renewcommand\ccode{#2\xspace}}
\newcommand\coursetitle[2][]{
  \ifthenelse{\isempty{#1}}%
    {}% do nothing if #1 is empty, else:
    {\renewcommand\cseries{#1\\[1ex]}}
  \renewcommand\ctitle{#2\xspace}}
\newcommand\semester[1]{
  \renewcommand\csemester{#1\xspace}}
\newcommand\meets[3]{
  \newcommand\AM{\textsc{am}}
  \newcommand\PM{\textsc{pm}}
  \renewcommand\cmeetson{#1\xspace}
  \renewcommand\cmeetsat{\Int{From 9:00 \AM}{#2}\xspace}
  \renewcommand\cmeetsin{\Int{\href{https://smu.brightspace.com/d2l/login}{Brightspace}}{#3}\xspace}}
\newcommand\website[2][]{
  \renewcommand\cwebsite{
  \ifthenelse{\isempty{#1}}%
    {\href{https://#2}}% copy url if #1 is empty, else
    {\href{#1}}% use specified url and print #2
  {#2}\xspace}}
\newcommand\cdescription[2][BF 1001; RM 1000 or GTRS 6000]{
  \renewcommand\cprereqs{#1}
  \renewcommand\cdescrip{#2\par}}
\newcommand\objectives[1]{
  \renewcommand\edobject{#1\par}}


\coursecode[7.0]{BSTH 1013}
\coursetitle{Old Testament Scripture}

\professor{Daniel R. Driver}
\phone{416.226.6620 x2201}
\email{ddriver@tyndale.ca}
\officehours{Mon. \& Thur., 1:30--3:30 \PM}

\semester{Fall 2014}
\meets{Tuesdays \& Thursdays} % \meets{on}{at}{in}
      {3:45--5:05 \PM}
      {Ballyconnor 2087}
\website{classes.tyndale.ca}
\cdescription{% copy from tyndale.ca/registrar/calendar
  Survey of the redemptive story in the three major divisions of the Old
  Testament (the Law/Torah, the Prophets, and the Writings), including
  an orientation to the historical backgrounds, religious context,
  literary forms, apocryphal dimensions, prophetic elements, matters of
  canon, text, interpretation, and critical issues.

  {\itshape Note that this class is both a general studies requirement
  and a prerequisite for advanced courses in Biblical Studies \&
  Theology.}
}

\ProvidesFile{preamble.tex}[2014/09/02 v1.1 -- Syllabus preamble]

% basic typography
\usepackage{fontspec,realscripts}
\defaultfontfeatures{Ligatures=TeX}
\setmainfont[%
  Numbers={Proportional,OldStyle},
  SmallCapsFeatures={RawFeature=+c2sc;+smcp,LetterSpace=1}]{Scala OT}
\newfontfamily\lining[Numbers={Lining}]{Scala OT}
\setsansfont[%
  SmallCapsFeatures={RawFeature=+c2sc;+smcp,LetterSpace=1}]{Scala Sans Pro}
\setmonofont[Scale=MatchLowercase]{Menlo}
\usepackage{sectsty}
\allsectionsfont{\sffamily\scshape}
\frenchspacing
\setlength{\emergencystretch}{3em} % prevent overfull lines

% custom font size and leading
\renewcommand\tiny{\fontsize{6}{9}\selectfont}
\renewcommand\scriptsize{\fontsize{7}{10}\selectfont}
\renewcommand\footnotesize{\fontsize{8}{10}\selectfont}
\renewcommand\small{\fontsize{9}{11}\selectfont}
\renewcommand\normalsize{\fontsize{10}{12}\selectfont}% base size
\renewcommand\large{\fontsize{11}{13}\selectfont}
\renewcommand\Large{\fontsize{12}{14}\selectfont}
\renewcommand\LARGE{\fontsize{16}{18}\selectfont}% "course syllabus \\ semester" benefits from more lead
\renewcommand\huge{\fontsize{20}{22}\selectfont}
\renewcommand\Huge{\fontsize{24}{26}\selectfont}

% layout packages: page, logo, tables
\usepackage[scale={0.6,0.8},
            xetex]{geometry}
\usepackage{graphicx}
\usepackage{array}     % allow insertions of column styling with >{}
\usepackage{booktabs}  % elegant horizontal rules in tables

% custom macros for a session count in the schedule of readings
\newcounter{session}
\newcounter{columns}
\newcounter{spancolumns}
\newcounter{midcolumns}
\newcounter{courseunit}
\newcommand\setcolumncount[2][0]{ % optionally set count to other than 0,
  \setcounter{session}{#1}        % e.g. to -1, or to a standing count
  \setcounter{columns}{#2}
  \setcounter{spancolumns}{#2-1}
  \setcounter{midcolumns}{#2-2}}
\newcommand\sessioncount{\stepcounter{session}\arabic{session}}
\newcommand\sessionskip[1]{\multicolumn{1}{@{}r@{ }}{#1}}
\newcommand\unit[1]{\multicolumn{\thecolumns}{c}{%
  \scshape\stepcounter{courseunit}\roman{courseunit}. #1}}
\newcommand\noclass[1]{\sessionskip{---.}&\multicolumn{\themidcolumns}{@{}l}{\itshape No Class: #1}}
\newcommand\rarr{\Large\char"EA31\hspace*{0.5pt}} % pointer finger w/ reminder tie in Scala OT
\newcommand\reminder[2]{\sessionskip{\rarr}&\multicolumn{\thespancolumns}{@{}l}{#1 \hfill #2}}

% color to match Tyndale's branding
\usepackage[usenames]{xcolor}
% predefined: black, white, red, green, blue, cyan, magenta, yellow
\definecolor{TyndaleURLs}{HTML}{0062A0} % links on tyndale.ca
\definecolor{TyndaleBlue}{cmyk}{1,1,0,.32}
\definecolor{TyndaleGold}{cmyk}{0,.27,1,0}
\definecolor{TyndaleRed}{cmyk}{0,1,.99,.04}
\definecolor{TyndaleBlack}{cmyk}{0,0,0,1}
\definecolor{TyndaleGreen}{cmyk}{.45,0,1,.24}
\definecolor{TyndaleOrange}{cmyk}{0,.79,1,0}
\definecolor{TyndaleAqua}{cmyk}{.47,0,.24,0}
\definecolor{TyndaleYellow}{cmyk}{.03,.03,.35,0}

% metadata (assumes a host of definitions are made in the main file)
\usepackage[setpagesize=false,     % leave this to geometry
            hyperfootnotes=false,  % fragile and distracting
            xetex]{hyperref}
\hypersetup{breaklinks=true,       % allow link text to break across lines
            colorlinks=true,       % colorlinks resets pdfborder to 0 0 0
            urlcolor=TyndaleURLs,  % for external links
            linkcolor=TyndaleRed,  % for normal internal links
            citecolor=TyndaleGold, % for bibliographical citations in text
            pdfauthor={\prof},
            pdftitle={\ccode: \ctitle},
            pdfsubject={Tyndale UC, \csemester},
            pdfcreator={github.com/danieldriver/syllabus}}
\urlstyle{same}                    % don't use monospace font for urls

% custom footlines
\usepackage{fancyhdr}
\pagestyle{fancy} % turn it on
\fancyhf{}        % reset everything
\renewcommand{\headrulewidth}{0pt} % remove header line as well
\lfoot{\sffamily\scshape\footnotesize\ctitle, \csemester}
\rfoot{\sffamily\scshape\footnotesize\prof\quad\thepage}

% language support: polyglossia needs to go last
\usepackage{polyglossia}
\setdefaultlanguage{english}
\setotherlanguage{hebrew}
\newfontfamily\hebrewfont[Script=Hebrew]{SBL BibLit}

% protect positioning of margin table in policy/grades -- needs to come after polyglossia
\usepackage{marginfix}

% gratuitous with custom title page, but useful as a fallback
\title{\ccode: \ctitle}
\author{\professor}
\date{\semester}


\begin{document}
\ProvidesFile{title.tex}[2013/09/06 v1.0 -- Syllabus title page]

\begin{titlepage}
  \begin{center}

    \begin{minipage}{\textwidth}
      \parbox[t]{0.5\textwidth}{
        \mbox{}\\[-\baselineskip] % dummy line to align parboxes
        \includegraphics[width=0.5\textwidth]{.syllabus/includes/TyndaleUC}}
      \hfill
      \parbox[t]{0.4\textwidth}{
        \raggedleft\LARGE\sffamily Course Syllabus\\
        \semester}
    \end{minipage}

    \vfill

    {\sffamily\textsc{\LARGE\MakeLowercase\coursecode}\\[2ex]
      \bfseries\itshape\Huge\coursetitle}

    \vfill

    \begin{tabular}{>{\bfseries}rl>{\bfseries}rl}
      \toprule
      Instructor & \professor                   & Course  & \liningnums{\coursecode}~\texttt{(v\version)} \\
      \midrule
      Phone      & \texttt{\phone}              & Meets   & \meeton                                       \\
      Email      & \href{mailto:\email}{\email} & Time    & \meetat                                       \\
      Office     & \officenumber                & Room    & \meetin                                       \\
      Hours      & \officehours                 & Website & \href{http://\website}{\website}              \\
      \bottomrule
    \end{tabular}

    \vfill

    \begin{description}
      \item[Commuter Hotline]
        Class cancellations due to inclement weather or illness will
        be announced on the commuter hotline at \texttt{416.226.6620
        x2187}. Alternately, weather cancellation information is posted
        at \href{http://tyndale.ca/weather}{tyndale.ca/weather}.
      \item[MyTyndale.ca]
        This course may have materials stored on its website, such as
        handouts or readings that may be needed in order to complete
        assignments. Students are responsible for checking these course
        pages on a regular basis. Here, too, students are able to view
        their grades throughout the semester. For more information see
        Section~\ref{mytyndale}, below.
      \item[Mail]
        Students are responsible for information communicated through
        their campus mailboxes and student e-mail accounts. A mailbox
        directory hangs beside the mailboxes. For more information
        contact the Registrar's office.
    \end{description}

  \end{center}

  \section{Course Description}
  \label{description}

  \emph{From the Academic Calendar:} \cdescription

\end{titlepage}
\setcounter{page}{2} % count the title page as page 1


\section{Learning Outcomes}
\label{outcomes}

The learning outcomes in this course will be determined by each
student's ability to meet specific goals (learning objectives). By the
end of the course students should be able to:
  name major Old Testament people and events;
  identify representative biblical quotes by book and chapter;
  give key dates for Israel's history and summarize the succession of
    superpowers in the Ancient Near Eastern political theatre from the
    time of Egypt to Greece;
%  contrast inner- and extra-biblical dating schemes;
  locate ancient Israel on a map;
  understand the general shape of the Masoretic Text tradition and
    differentiate it from other canonical orders and traditions;
  recognize genres of biblical literature and cite examples from the
    reading;
  list the Ten Commandments in order;
  articulate multiple rationales for sabbath observance;
  defend a decision to pronounce or circumlocute the divine name;
  classify prophetic literature relative to the exile;
%  report on instances of duplication in both the Law and the Prophets;
  memorize and recite Psalm 8.

Students should also be able to identify settings in which the
scriptures of Israel are read (notably the synagogue, church and
academy), employ terminology appropriate to these reading communities,
recognize where their own biographies and commitments place them in
relation to the Hebrew Bible/Old Testament and its uses, and monitor and
test their individual attitudes and assumptions. They should be able
to extend their awareness of the Bible's reading communities to the
Bible's long history of reception. Finally, students should be able to
infer what Jesus might have meant in speaking of ``the law of Moses, the
prophets, and the psalms'' (Luke 24:44), and begin to hear claims about
New Testament fulfilment of scripture in light of the unique voice that
the Old Testament retains along side of the New in Christian scripture.

\section{Required Texts \& Materials}
\label{texts}

All required textbooks are available in the Tyndale Bookstore. The
course pack is available in class, directly from Dr. Driver.

\begingroup
\renewcommand{\section}[2]{}% temporarily remove the section heading
\begin{thebibliography}{Heschel}% use the longest item in the bibliography
  \bibitem[NRSV]{nrsv} \emph{The New Oxford Annotated Bible with Apocrypha}.
    Oxford: Oxford University Press, 2010 (4\textsuperscript{th}
    edition). Many serviceable translations of the Bible exist, but I
    assign the NRSV because it is a standard among biblical scholars. If
    you do not already have a copy then you must purchase one for this
    class.
  \bibitem[Heschel]{heschel} Abraham Joshua Heschel.
    \emph{The Sabbath: Its Meaning for Modern Man}. New York: Farrar,
    Straus and Giroux, 1951 (repr. 2005). Read this classic example of
    20\textsuperscript{th} century Jewish spirituality on your own, no
    later than the fourth week of class. Review it in preparation for the
    class on which it is assigned.
  \bibitem[Moberly]{Moberly} R.\,W.\,L. Moberly.
    \emph{Old Testament Theology: Reading the Hebrew Bible as Christian
    Scripture}. Grand Rapids: Baker Academic, 2013. Read these eight
    essays by a leading Old Testament scholar and theologian in parallel
    with the assigned Bible readings. It will be best if you read the
    primary texts first, of course.
  \bibitem[Driver]{driver} Daniel R. Driver.
    \emph{\ctitle\ Introductory Course Pack}. Toronto, \csemester.
    In response to the requests of many students, I have compiled a pack
    of notes for this class. I will circulate the relevant pages before
    each class. It behooves you to gather, annotate, and study this
    material carefully.
\end{thebibliography}
\endgroup

\section{Supplementary Texts}
\label{supplementary}

Supplementary readings are listed in the bibliographies that follow each
page of notes in the course pack. Excerpts from this literature,
ordinarily an article or a book chapter per week, will be available for
download through the course website (alternately, some may be kept on
reserve in the library). Students are not strictly required to read this
supplementary material; then again, students who choose not to read it
should not expect to get an A for the course.

Note that the required readings for this survey emphasize a selection of
primary texts from the Old Testament. Reading in the vast body of
secondary literature on the Bible is largely left for upper-division
courses. Diligent students, however, and all prospective majors, will
make a special point of reviewing the supplementary material carefully.

\section{Evaluation}
\label{evaluation}

\subsection{Grade Structure for \ccode}
\label{structure}

\begin{enumerate}
 \item At various points throughout the semester I will announce
   \textbf{reading quizzes}. They are designed to ensure that you have
   read the assigned material carefully, and may be administered online
   (outside of class) or on paper (in class). Online quizzes are open
   book, but they are timed, and only a single attempt is allowed.
   Quizzes given in class may not be made up in the case of absence.
 \item Students are to \textbf{memorize Psalm 8} in the KJV and recite
   it to me privately, in my office, in the fourth or fifth week of
   class. A sign-up sheet will be circulated in week three.
 \item A \textbf{midterm exam} will be given in class as scheduled in
   the course outline (Section~\ref{outline}). At the instructor's
   option, it may include a take-home writing component.
 \item A comprehensive \textbf{final exam} will be held during the exam
   period as schedule by the Registrar (see Section~\ref{exams}). It is
   designed to assess learning outcomes---how well you have obtained the
   objectives outlined in Section~\ref{outcomes}. Test your readiness
   for the final by asking yourself how well you can perform the tasks
   outlined in those paragraphs.
\end{enumerate}

The breakdown for the semester's total work is as follows:

\begin{table}[htbp]
  \centering
  \begin{tabular}{lr}
    \toprule
    Reading Quizzes & 25\% \\
    Memorization    & 15\% \\
    Midterm Exam    & 30\% \\
    Final Exam      & 30\% \\
    \bottomrule
  \end{tabular}
  \caption{Distribution of Grades}
  \label{distribution}
\end{table}

\ProvidesFile{grades.tex}[2023/09/06 v2.5 -- Course policy]

AST's \href{https://www.astheology.ns.ca/students/resources.html}{Academic
Calendar} provides guidelines and criteria for academic assessment. Marks are
assigned by letter grade using these benchmarks. Note that graduate and
undergraduate programs are currently assessed on a different scales at AST (see
\autoref{grade-syst}). For graduate students, a grade of C is deemed
``Unsatisfactory,'' whereas for undergraduate (BTh) students, a C is
``Satisfactory'' and the grade of D is possible for ``Marginal'' performance.
Other grades have different definitions, too.

\begin{table}[htbp]
 \centering
 {\lining\footnotesize
 \begin{tabular}{lll|ll}
   \toprule
               & \multicolumn{2}{c|}{Graduate Courses}   & \multicolumn{2}{c}{Undergraduate Courses} \\
               \cline{2-5}
   A+          & 94--100    & Exceptional               & 90--100    & Excellent                 \\
   A           & 87--93     & Outstanding               & 85--89     & Excellent                 \\
   A\char"2212 & 80--86     & Excellent                 & 80--84     & Excellent                 \\ [1ex]
   B+          & 77--79     & Good                      & 75--79     & Good                      \\
   B           & 73--76     & Acceptable                & 70--74     & Good                      \\
   B\char"2212 & 70--72     & Marginally Acceptable     &            &                           \\ [1ex]
   C           & 60--69     & Unsatisfactory            & 60--69     & Satisfactory              \\
   D           &            &                           & 50--59     & Marginal                  \\
   F           & 0--59      & Failure                   & 0--49      & Failure                   \\
   FP          & 0          & Failure due to Plagiarism & 0          & Failure due to Plagiarism \\
   \bottomrule
 \end{tabular}}
 \caption{Outline of Grades at AST}
 \label{grade-syst}
\end{table}


\section{Policy on Assignments \& Exams}
\label{policy}

\ProvidesFile{academic_calendar.tex}[2015/09/10 v1.1 -- Course policy]

\newcommand{\AC}{Academic Calendar}
\newcommand{\SecAC}{the \AC}% Starting in 2015, the AC no longer has sections!

All policy in Sections~\ref{policy}, \ref{expectations} and \ref{support} of
this syllabus applies to this course in addition to policy in the current
\href{http://www.tyndale.ca/registrar/calendar}{\AC}. In some cases the
syllabus underscores the general policy, while in other cases it supersedes it.

For all matters not covered in this syllabus, refer to \SecAC, ``University
College Academic Policies, Procedures, and Notices.'' Students are strongly
encouraged to read this document carefully at least once in their career at
Tyndale, and to review it every year they matriculate.

\ProvidesFile{assignments.tex}[2014/09/02 v1.1 -- Course policy]

\subsection{Assignments}
\label{assignments}

This is a university course. All papers and other writing assignments
should therefore be written at the university level. Submissions must be
typewritten and double-spaced, should be free from error, and in this course
should follow the \emph{SBL Handbook of Style} (refer to the free, online
\href{https://www.sbl-site.org/assets/pdfs/pubs/SBLHSsupp2015-02.pdf}{SBLHS
Student Supplement}.)

If you ever struggle with composition---anything from the relatively simple
matters of spelling, grammar and proper citation to deeper-level issues of
tone, structure and argument---then please make use of the Writing Centre (see
Section~\ref{writing_centre}). Experienced writers know that drafts and peer
feedback are integral to the writing process. Inexperienced writers are often
unaware that their surface-level errors create credibility problems with their
readers. When you \href{http://theoatmeal.com/comics/misspelling}{misspell
common words}, fail to know \href{http://theoatmeal.com/comics/apostrophe}{how
to use an apostrophe}, or do not bother to cite your sources correctly, why
should your readers trust you with more important matters like the facts and
ideas under discussion?

\subsubsection{Deadlines}
\label{deadlines}

Assignments \emph{must} be submitted on time. Even if the work is rough or
incomplete, you must turn in something by the due date to receive any credit
whatsoever. Unless I specify differently in class, papers and take-home exams
are due by 11:59 \textsc{pm} on the due date. All other work is due at the
start of the day's class.

Note that, because no late work is accepted in this class, there is no scale
of penalty for unexcused late assignments. If a truly extraordinary event
keeps you from doing your best work, then let me know so that we can make
special arrangements. I am guided by the \AC\ in what counts as extenuation.
``Extensions are not granted for what best could be described as `poor time
management' or `over-involvement' in an extracurricular activity.''

\subsubsection{Submission as PDFs}
\label{submission}

Papers and some other assignments in this course are to be submitted
electronically through the course pages (Section~\ref{lms}). To preserve
formatting, formal writing assignments must be uploaded in Portable Document
Format. There are many ways of creating PDFs; it is your responsibility to
know how to do so on the computer platform you use, and to generate and submit
your PDFs on time.

\subsubsection{Backup}
\label{backup}

In the event of the loss of assignments post-submission---electronic systems
fail, and my office has flooded before---students are required to keep backup
copies of all assignments submitted.

Learning how to secure and preserve your work is a peculiar challenge of the
digital age. Plan on the crash of your hard drive, and the theft of your
laptop (the first is inevitable, the second quite probable). If you do not
have a backup strategy, I recommend that you start with a free account on
\href{http://db.tt/U7eP1vs}{dropbox.com}.

\ProvidesFile{exams.tex}[2013/08/19 v1.0 -- Course policy]

\subsection{Examinations}
\label{exams}

My examination policy follows that outlined in \SecAC, part of which is
summarized below for emphasis.

\begin{enumerate}
  \item
    Midterm exams will be held as scheduled by the instructor. If you miss
    the exam for a legitimate reason, you must write the exam within the same
    number of days that you were absent from school (possibly the next day).
  \item
    Final examinations will take place during the exam period as scheduled
    by the Registrar. Students are responsible for noting the date, time and
    location of their final exam in this class. Students are also responsible
    for familiarizing themselves with the Registrar's examination policies.
  \item
    The following rules apply to every final examination:
    \begin{enumerate}
      \item
        No student is permitted to take into the examination room any materials
        relating to the examination subject, including Bibles.
      \item
        No student may leave the room without permission from the proctor.
      \item
        No student may leave his or her seat during the final fifteen minutes.
      \item
        Students must not linger in the halls outside the examination rooms
        while examinations are being written.
      \item
        No student will be permitted to write beyond the allotted time without
        special permission of the Registrar (see Section~\ref{accommodation}).
    \end{enumerate}
  \item
    Provisions exist for students who are justifiably unable to write the final
    exam at the scheduled time. See the \AC\ for details, and make arrangements
    through the Office of the Registrar.
  \item
    Normally, a final exam can only be reschedule in two circumstances:
    (a) a documented illness, or (b) a conflict with another
    exam (two at the same time, or three within 24 hours).
    \href{http://www.tyndale.ca/registrar/final-exam-schedule-and-policies}
    {Apply to the Registrar} in either case.
\end{enumerate}


\section{Student Expectations \& Guidelines}
\label{expectations}

\ProvidesFile{academic_integrity.tex}[2013/08/19 v1.0 -- Course policy]

\subsection{Academic Integrity}
\label{integrity}

Integrity in academic work is required of all students. Academic dishonesty
is any breach of this integrity. It includes such practices as cheating (the
use of unauthorized material on tests and examinations), submitting the same
work for different classes without permission of the instructors, using false
information in an assignment (including false references to secondary sources),
improper or unacknowledged collaboration with other students, and plagiarism.

Tyndale takes seriously its responsibility to uphold academic integrity,
and to apply consequences for academic dishonesty. Consult \SecAC\ for more
information on the school's policy and its application to your work in this
course.

\ProvidesFile{attendance.tex}[2013/08/22 v1.0 -- Course policy]

\subsection{Attendance}
\label{attendance}

``Faithful attendance at classes is an important indicator of student maturity
and involvement'' (\AC). Remember, too, that you are responsible for everything
that happens in every class. Your best policy is to attend and engage. Please
do not ask me to repeat for your benefit anything I have said in a class you
have missed.

Keeping a record of attendance is mandatory for faculty at Tyndale (in contrast
to many other colleges and universities). The University College publishes
guidelines for how attendance should bear on your final evaluation in a course,
and I adhere to them. Note that four lates equals one absence.

What should you do if you miss an undue number of classes? First, arrange for
a classmate to brief you on the material missed, or get my permission for a
classmate to make a recording for you (see Section~\ref{recording}). Second,
notify the Dean of Students in person or by phone. If illness is the cause you
will need to submit a doctor's certificate upon return. The Dean of Students
will notify your professors of the reason for the absence and suggest that they
take this into consideration when assigning grades.

\ProvidesFile{technology.tex}[2014/09/02 v1.1 -- Course policy]

\subsection{Technology}
\label{technology}

Technological innovation has brought students and educators a number of
powerful new tools, and I encourage you to use them as you research, write, and
collaborate. Some of these tools also call for disciplined use and management.

\subsubsection{Email}
\label{email}

Email can be a chore, and you may prefer other channels of communication. As a
matter of policy, however, students must use their myTyndale accounts for all
course-related email correspondence. During term time you should check your
school account at least once a day (optional on weekends). I myself aim to
check my school email at the beginning and end of each workday. At other times
my email client is often closed. I will try to answer your messages within 24
hours, though you should not expect replies on weekends.

Try to keep your messages topical and brief. I would vastly prefer to conduct
any conversations of substance in person, or else over the phone. Please note
and make use of my office hours. I would happily receive an email from you
requesting an alternate meeting time if my regular hours do not suit your
schedule.

% Note the prof who adopted an anti-email policy! https://www.insidehighered.com/news/2014/08/27/sake-student-faculty-interaction-professor-bans-student-email?utm_source=slate&utm_medium=referral&utm_term=partner

\subsubsection{MyTyndale.ca / classes.tyndale.ca}
\label{mytyndale}

Tyndale's course pages are an efficient means of distributing articles, notes,
slides, and other course-related materials. This is also where instructors log
attendance and upload grades for assignments. Students are therefore required to
check the site for updates about their classes as well as for any materials
needed for lectures and assignments.

My own use of this platform varies from semester to semester, and from
course to course. At times I may ask you to use the forums, quiz module, or
other parts of the system. At a minimum I will use the site as a repository
for course materials, and as a destination for your submission of PDFs
(Section~\ref{submission}).

\subsubsection{Laptops and Other Devices}
\label{laptops}

Use of laptops is forbidden in my classroom, except to facilitate presentation.
I implement this policy because of the cognitive costs of multitasking, with the
aim of giving you and your peers the best chance of success. I also hope to
foster a culture of keen intellectual engagement.

As \href{http://dx.doi.org/10.1016/j.compedu.2012.10.003}{cognitive
psychologists at McMaster and York Universities demonstrated in 2013}, ``laptop
multitasking hinders classroom learning for both users and nearby peers.'' There
is little new in their finding that the allegedly multitasking student does less
well in class (11\% worse on the quiz in their experiment). This effect has been
shown many times. Rather, their novel discovery is that classmates
\emph{without} laptops who sat with a \emph{view} of another student's screen
did worse than the students who had a computer (17\% worse than those with no
laptop in sight).

Prohibiting laptops is not the only possible response to these findings.
However, there is evidence that \href{http://on.wsj.com/pjtJaK}{writing by hand}
brings a number of cognitive benefits, and a
\href{http://dx.doi.org/10.1177/0956797614524581}{2014 Princeton University
study} ``found that students who took notes on laptops performed worse on
conceptual questions than students who took notes longhand.'' If you are a heavy
laptop user then consider this an opportunity to experiment with different
technologies in the classroom.

As for the myriad networked devices that many of us carry, it's a simple matter
of professionalism to keep these things silent and out of sight. E-readers and
tablets are permitted \emph{only if they are used to display the assigned
reading}. If this is how you choose to read, let me invite you to put the
machine in airplane mode while class is in session.

\subsubsection{Recording of Classes}
\label{recording}

Students must request permission from the professor of any class that they
would like to record. Where permission is granted, students are expected to
supply their own equipment. In general I prefer \emph{not} to have my classes
recorded, and I am not at all friendly to being recorded without my knowledge.
In cases where I grant permission, I stipulate that the recordings must be for
personal use only. They should not be shared with other students, even with
students in the same section, and they absolutely must not be posted online or
otherwise distributed.

If a student is not able to attend a lecture and would like to have it recorded,
it is the responsibility of the student first to obtain the professor's
permission, and then to find another student to record the lecture. I will not
make the recording for you.

\ProvidesFile{support.tex}[2014/08/06 v1.1 -- Course policy]

\section{Student Support}
\label{support}

\subsection{Tyndale Writing Centre}
\label{writing_centre}

Through a combination of tutorials, workshops and resources, Tyndale’s Writing
Centre offers a comprehensive program of writing support to Tyndale students,
including individual 30-minute tutoring sessions. Students may bring essays that
have been graded (and, at least for my classes, essays that have not yet been
submitted for a grade) and will receive detailed suggestions for improving their
writing. This service, at no charge to students, is available by appointment.

Professors may recommend that a student go to the Writing Centre for help:
students are strongly encouraged to follow such recommendations. The Academic
Standards Committee may require an undergraduate student who is experiencing
difficulty in his or her academic program to go to the Writing Centre for
assistance and support. Many top students also elect to go.

\subsection{Tyndale UC Tutoring Program}
\label{tutoring}

Tyndale University College is committed to helping its students achieve academic
success. For this purpose, students in need of academic assistance may request
peer tutoring, free of charge, in each academic department. This includes
students on academic probation, students who have received failing grades in a
course or courses, or students who have been referred for tutoring by their
instructor.

For more information on scheduling tutoring appointments, or for those
interested in becoming peer tutors, students may contact the Office of the
Senior Vice President Academic or their respective University College department
chairs.

\subsection{Accommodation}
\label{accommodation}

Students with documented disabilities may be granted special accommodation for
exams, and in some cases for other assignments. It is even possible to get
permission to use a laptop in class (Section~\ref{laptops}), although I will
need to be convinced of the use case. It is up to the student to contact the
Dean of Students as early as possible in the semester---not later than the
second week---and to document the need. The Dean of Students will then advise
each of the student's professors of the accommodations that may be required.
Please note that special arrangements for assignments need to be made with me
well in advance of assignment due dates (Section~\ref{deadlines}). Timely
requests shall not unreasonably be denied.


\section{Course Outline}
\label{outline}

Readings of a given day should be completed before the start of class on
the same date. Primary readings, Heschel and Moberly are compulsory. Any
supplementary readings that I circulate (per
Section~\ref{supplementary}) will be keyed to the session number (listed
under \textbf{\S} in Table~\ref{schedule}), and should ideally be read
before class as well. We will adhere to the schedule as closely as
possible, though I reserve the right to adjust it to suit the needs of
the class.

\newcommand\Yhwh{\textsc{Yhwh}}
\newcommand\rarr{\char"2192\hspace*{0.5pt}}

\setcolumncount{4} % set up \sessioncount, \unit{} and \noclass{} macros
\begin{table}[phtb]
  \centering
  \begin{tabular}{>{\sessioncount.}r@{ }llr}
    \toprule
    \sessionskip{\textbf{\S}.}&\textbf{Primary Reading}&\textbf{Secondary Reading}&\textbf{Date}\\
    \midrule

    \unit{Beginnings} \\

          & Psalms 8, 74, 104               & Syllabus         &  4 Sep.     \\
          
          & Genesis 1--11                   & Heschel (begin)  &  9 Sep.     \\
          & Genesis 12--36                  &                  & 11 Sep.     \\
          
          & Genesis 37--50                  &                  & 16 Sep.     \\
          & Exodus 1--6, 32--34             &                  & 18 Sep.     \\ [1ex]

    \unit{Out of Egypt} \\
          
          & Exodus 7--24                    & Moberly 3        & 23 Sep.     \\
          & Exodus 20, Deuteronomy 5        & Heschel (finish) & 25 Sep.     \\
          
          & Deuteronomy 1--11               & Moberly 1        & 30 Sep.     \\
          & Deuteronomy 12--26              &                  &  2 Oct.     \\
          
          & Deuteronomy 27--34              & Moberly 2        &  7 Oct.     \\ [1ex]

    \unit{In the Land} \\

          & Joshua 1--13:7, 22--24          &                  &  9 Oct.     \\
          
          & Judges 1--16, 19--21            & Moberly 4        & 14 Oct.     \\
          & 1 Samuel 1--15                  &                  & 16 Oct.     \\
          
          & \textbf{Midterm Exam} (\& Essay Due If Assigned) & & 21 Oct.     \\
    \noclass{Reading Days}                                     & 23--24 Oct. \\
    
          & 1 Samuel 16--31, 2 Samuel 1--12 &                  & 28 Oct.     \\
          & 1 Kings 1--14                   &                  & 30 Oct.     \\
          
          & 2 Kings 17--25                  &                  &  4 Nov.     \\ [1ex]
    \unit{Latter Prophets} \\
          & Isaiah 1--12, 36--39            & Moberly 5        &  6 Nov.     \\

          & Hosea 1-14, Joel 1--3           &                  & 11 Nov.     \\
          & Amos 1--9, Obadiah, Jonah 1--4  & Moberly 6        & 13 Nov.     \\ [1ex]
          
    \unit{Wisdom \& Praise} \\

          & Job 1--11                       &                  & 18 Nov.     \\
          & Job 27--31, 38--42              & Moberly 8        & 20 Nov.     \\ 
          
          & Ecclesiastes 1--12              &                  & 25 Nov.     \\
          & Psalms 1--2, 44, 51, 72--73, 89--90, 105--106, 145, 150 & Moberly 7 & 27 Nov. \\ [1ex]
          
    \noclass{Reading Days}                                               & 2--3 Dec.   \\
    \bottomrule
  \end{tabular}
  \caption{Schedule of Readings}
  \label{schedule}
\end{table}

See the \href{http://www.tyndale.ca/registrar/important-dates}{%
Registrar's website} for a list of other important dates. Generally, the
last day to add or drop a class without penalty is the end of the second
week of class.

\section{Course Bibliography}
\label{bibliography}

An introductory bibliography is provided at the end of each page of
notes in the course pack. Pursue this material as your interest
dictates, or if you wish to check my handling of any of the material I
present. Pay special attention to items flagged as supplementary reading
(Section~\ref{supplementary}).

\end{document}
